\subNsection{1}{Endliche Automaten}{Kristian Popov} \label{4_2_EA}

Im Thema \emph{Endliche Automaten} können Studierende deterministische und
nichtdeterministische endliche Automaten selbst definieren und Eingabewörter Schritt
für Schritt ausführen.
Jeder Simulationsschritt wird gleichzeitig in der Übergangstabelle und im
Zustandsdiagramm dargestellt.
Ergänzt wird die Einzelsimulation durch einen Stapeltest für mehrere Wörter, die
Umwandlung rechtslinearer Grammatiken in Automaten und Übungen mit automatischer
Rückmeldung.

\subNsection{2}{Entscheidungen in der Architektur}{Kristian Popov}

Ursprünglich war nur die Unterstützung deterministischer Automaten vorgesehen.
Nichtdeterministische kamen hinzu, weil dort ein Symbol in mehrere Zustände zugleich
führen darf: Ein Automat ist so oft direkt aus der Sprachbeschreibung ablesbar und
kompakter, und die Mehrdeutigkeit muss erst beim Schritt zum DFA aufgelöst werden.
Da jeder DFA gleichzeitig ein NFA ist, wurde für beide Varianten bewusst nur eine
Ausführungslogik verwendet.
Der Kern arbeitet mit einer Menge aktiver Zustände, die bei einem DFA höchstens ein
Element enthält; $\varepsilon$-Übergänge werden dabei stets über die
$\varepsilon$-Hülle aufgelöst.
Der deklarierte Automatentyp steuert Validierung und Eingabemaske, nicht aber die
Ausführung.

Der Simulator ist in fünf Bereiche aufgeteilt:

\begin{description}[leftmargin=*,style=nextline,itemsep=0.75ex]
\item[\texttt{domain}]
enthält die Fachlogik wie Simulation, Validierung, Grammatikumwandlung, Stapeltest
und Übungskontrolle.

\item[\texttt{application}]
verbindet die Fachlogik mit React und verwaltet Simulation und Editor.

\item[\texttt{adapters}]
bereitet die Daten für die Übergangstabelle auf.

\item[\texttt{presentation}]
enthält die Benutzeroberfläche mit Diagramm, Tabelle, Eingabewort, Stapeltest und
Grammatikeingabe.

\item[\texttt{config}]
enthält Standardautomaten und Beispiele.
\end{description}

Die Trennung sorgt dafür, dass die Fachlogik unabhängig von React bleibt und ohne
gerenderte Oberfläche testbar ist.
Das Thema \emph{Reguläre Sprachen} folgt demselben Schnitt und greift zusätzlich
direkt auf die Grammatikumwandlung dieses Themas zurück.

Ein Simulationsschritt erzeugt einen neuen Zustand mit der bisherigen Historie.
Dadurch sind sowohl das Zurückgehen als auch die Auswahl früherer Schritte möglich,
ohne eine zusätzliche Historie in der Oberfläche zu verwalten.

Auch die Validierung ist von der Oberfläche getrennt.
Sie liefert Fehlercodes und relevante Werte, während die Oberfläche daraus die
passenden Meldungen erzeugt.
Dadurch bleibt die Fachlogik unabhängig von der Sprache der Benutzeroberfläche.

Die mitgelieferten Automaten sind bewusst partiell definiert.
Fehlt ein Übergang, wird das Wort abgelehnt.
Auf einen zusätzlichen Fangzustand wurde verzichtet, da dieser zwar die Sprache nicht
verändert, aber die Diagramme unnötig vergrößert.

\subNsection{2}{Entscheidungen im UI}{Kristian Popov}

Ursprünglich war für die Simulation nur die Übergangstabelle vorgesehen.
Während der Entwicklung zeigte sich jedoch, dass der Ablauf eines Wortes damit schwer
nachzuvollziehen ist: Zyklen und Sackgassen werden erst sichtbar, wenn man mehrere
Zeilen gedanklich verknüpft.
Deshalb wurde zusätzlich ein Zustandsdiagramm integriert.

Beide Darstellungen sind an denselben Simulationsschritt gebunden: Die Tabelle hebt
genau die Zelle hervor, die dem im Diagramm markierten Übergang entspricht.
Sie zeigt die formale Definition des Automaten, während das Diagramm den
tatsächlichen Weg des Wortes sichtbar macht.
Bei größeren Automaten wird das automatisch berechnete Layout schnell
unübersichtlich.
Einzelne Zustände lassen sich deshalb frei ziehen, der Ausschnitt vergrößern und
verkleinern und die Ausgangsansicht mit einem Klick wiederherstellen.
Aktive und fehlerhafte Übergänge bleiben dabei auch bei geringer Zoomstufe
beschriftet.

Verstöße gegen den Determinismus werden nicht nur aufgelistet, sondern unmittelbar an
Zustand und Übergang im Diagramm markiert.
Damit ist die Ursache dort sichtbar, wo sie entsteht.

Der Editor umfasst alle wesentlichen Bestandteile eines Automaten: Automatentyp,
Alphabet, Zustände, Start- und Endzustände sowie das Eingabewort.
Die Zustandsmenge wird direkt über die Übergangstabelle verwaltet.
Alphabetsymbole sind auf einzelne Zeichen beschränkt, damit eine Eingabe wie
\verb|10| nicht versehentlich als ein einzelnes Symbol aufgenommen wird.
Die Zeile für $\varepsilon$-Übergänge erscheint nur bei einem NFA, sobald sie belegt
ist oder ausdrücklich angefordert wird.

\subNsection{2}{Andere Entscheidungen}{Kristian Popov}

Neben der Einzelsimulation gibt es einen Schrittverlauf, über den bereits erreichte
Schritte direkt ausgewählt werden können.
Der gewählte Schritt wird dabei vom Anfang des Laufs an neu berechnet statt
zwischengespeichert; bei den hier üblichen Wortlängen ist das unkritisch und hält
Anzeige und tatsächlichen Ablauf zwingend deckungsgleich.
Der Stapeltest ermöglicht es, mehrere Wörter gleichzeitig zu prüfen.
Ein Ergebnis kann anschließend direkt in die Einzelsimulation übernommen werden.

Eine weitere Funktion ist die Umwandlung rechtslinearer Grammatiken in Automaten.
Nichtterminale werden dabei zu Zuständen und Produktionen werden in entsprechende
Übergänge oder Endzustände übersetzt.
Für die Eingabe werden zwei Schreibweisen unterstützt, die automatisch erkannt
werden.

Beispiele und Übungen sind getrennt aufgebaut.
Beispiele enthalten fertige Automaten und dienen zum Nachvollziehen, während Übungen
mit einem leeren Automaten beginnen und zusätzlich Aufgabenstellung, Referenzlösung
und automatische Kontrolle enthalten.
Die Übungen sind nach \emph{Sprache} $\rightarrow$ \emph{Automat} und
\emph{Grammatik} $\rightarrow$ \emph{Automat} gegliedert.

Die Übungskontrolle vergleicht nicht die Struktur der Automaten, sondern deren
Verhalten.
Dazu werden alle Wörter über dem gemeinsamen Alphabet bis zu einer festgelegten Länge
geprüft.
Der Automat der Studierenden muss für jedes Wort dieselbe Entscheidung treffen wie
die Referenzlösung.
Bei einer Abweichung wird das erste unterschiedliche Wort als Gegenbeispiel
angezeigt.
Dadurch können auch unterschiedlich aufgebaute, aber funktional korrekte Automaten
akzeptiert werden.
Die Prüfung stellt jedoch keinen formalen Äquivalenznachweis dar.
Strukturelle Mängel wie ungültige Übergänge werden vorrangig gemeldet, da sie meist
die eigentliche Ursache sind und eine Liste fehlgeschlagener Testwörter dann in die
Irre führen würde.

Der Bearbeitungsstand jeder Übung wird gespeichert und beim erneuten Öffnen
wiederhergestellt, sodass der Wechsel zwischen Übungen keine bereits begonnene
Bearbeitung überschreibt.
